\scanpage{49}
% typed verso
% @top right, pagination
\mspage{6}
% @01–02, opening
Continuing in this way we construct the minimal model
% @03, first display
\[
 \Lambda(x_1,x_2,\ldots\mathbin{:}d)\underset{\rho}{\to}\mathcal{E}
\]
% @04, continuation
satisfying properties i) and ii) above.

% @05–06, theorem heading
\uline{Theorem B}: Let $X$ be a triangulated space which is simply connected. If
% @07, second display
\[
 \Lambda(x_1,x_2,\ldots)\underset{\rho}{\to}\mathcal{E}
\]
% @08, hypothesis
is a minimal model of the de Rham algebra then

% @09–10, conclusion i
i) the rank of $\pi_i X$ is the number of variables in degree $i$.

% @11–12, conclusion ii
ii) the rank of the Hurewicz homomorphism $\pi_i X\to H_i X$ is the number of $d$-closed generators in degree $i$.

% @13–15, conclusion iii
iii) the Whitehead products on homotopy are described by the quadratic terms of the $d$-formulae of generators.

% @16–23, proof discussion
Theorem B is proved by a not so standard Postnikov argument using Theorem A and the commutativity of the wedge product multiplication in $\mathcal{E}$. The argument is based on the beautifully simple Guy Hirsh method of computing the cohomology of a principal $G$-bundle over a manifold using differential forms (see appendix H). The Hirsch method was suggested by Rene Thom in the study of Postnikov systems in the Cartan Seminar 1954.

% @24–25, final paragraph
The proof of theorem B shows the model can be used to directly construct the \uline{rational} Postnikov system
% archive sheet number: 48

